\section{开元的中枢从清洗开始}
\label{sec:xuanzong-frontiers}

710年唐隆政变诛韦后集团，李隆基拥立睿宗；712年受禅，713年又除去太平公主集团。三次转折不能只被归结为年轻君主的决断：禁军、宗室、外朝官员和宫城空间都构成行动条件。正史把先后安排得很清楚，参与者的犹疑、命令传递和未参与者的处境却难以复原。可以确定的是，713年之后不再有两个并存的权力中心，玄宗得以更直接地调配中枢人事。%
\footnote{唐隆政变、玄宗即位与太平公主集团覆灭见 ST-710-TANGLONG-COUP、ST-712-XUANZONG-ACCESSION、ST-713-TAIPING，依据《旧唐书》《新唐书》睿宗、玄宗纪及《资治通鉴》。政变预谋和宫禁细节主要来自事后胜者叙事。}

开元九年宇文融括户检田、二十一年裴耀卿整顿江淮漕运，把所谓“开元之治”落到名册与水路上。前者试图把逃户、隐漏户重新放进赋役登记，后者以转运节点改善关中供粮；两者都显示朝廷需要不断修补人口、税源和粮道之间的断裂。诏令与传记能确认政策方向，不能让我们把括得的户数、漕运总量或地方负担当作统一可靠的全国统计。%
\footnote{括户检田与漕运整顿见 ST-721-YUWEN-RONG、ST-733-PEI-YAOQING，依据《旧唐书·食货志》《旧唐书》宇文融、裴耀卿传、《新唐书》相关志传及《通典》。制度条文和行政数字须分地区、口径核验。}

725年封禅泰山以礼仪展示政局与边功；736年李林甫拜相、737年三位皇子被废赐死而李亨继立，说明仪式的安定与继承的脆弱可以同时存在。李林甫的长期相位不能仅用一人品格解释，它也依靠文书、人事和皇帝授权；皇子事件同样不能从后来安史结局倒推为必然。天宝元年复郡名、改刺史为太守，是对行政语言的再次整齐化，而不保证地方财政或军政关系同步改变。%
\footnote{封禅、李林甫拜相、皇子废立与天宝改制见 ST-725-TAISHAN、ST-736-LI-LINFU、ST-737-CROWN-PRINCES、ST-742-TIANBAO，依据两《唐书》玄宗纪、李林甫传、礼仪和地理志及《资治通鉴》。人物归责、礼仪效果与地方执行皆须保留史料限度。}
